% ===== §14 =====
\section{The Emergence of Culture in Intimate Relation: A Definition}
\label{sec:definition}

\noindent
The synthesis can now state its definition. The thesis surveyed the traditions and mapped their common presupposition; the antithesis inverted that presupposition, gave the inversion its content in the exteriorization of drive, and secured it in a dialectic of the material and the ideal, all within the framework of the three-register dialectic of the relational Real. What these yield, gathered into a single formulation, is a definition of what culture is and how it emerges---general in its statement, and about to be made special, in the second part, to the intimate dyad from which the paper takes its name.

\subsection*{The definition}

\begin{claim}[Definition]
\emph{Culture}, in its mode of being, is generated as the dialectical exteriorization, by relational subjects and under the constraint of material conditions, of their interwoven structure of drive and their shared history of generation, into transmissible and re-enactable forms---idioms, rites, private references, understandings---across the three registers of the imaginary, the symbolic, and the real; a generation that, returning upon the relation that carries it, organizes and sustains that relation in turn.
\end{claim}

\noindent
Every clause of this definition carries a result established earlier, and the definition is best read by unfolding them. \emph{The dialectical exteriorization}: culture is spirit going out of itself into a visible form and returning altered, the Hegelian movement of \xref{sec:inversion}, not a deposit laid down once but a going-out that is a coming-to-self. \emph{By relational subjects}: the unit is neither the individual nor the collectivity but the relation, the coupling in which and for which the culture is made; this is the coordinate the map left vacant and the framework filled. \emph{Under the constraint of material conditions}: the exteriorization is materially constrained but not materially determined, bounded in its range by the forces and relations of production and open within that range, the dialectic of \xref{sec:material}. \emph{Of their interwoven structure of drive}: what is exteriorized is drive, coded and sublimated, and in intimate relation the drive of the two is interwoven rather than parallel, so that the culture is the sediment of a shared fantasy, the result of \xref{sec:drive}. \emph{And their shared history of generation}: culture is not only the exteriorization of a present structure but the sedimentation of a history, the accumulated record of what the relation has generated, folded into forms that can be taken up again. \emph{Into transmissible and re-enactable forms}: the sedimented side the traditions correctly saw is preserved here, as the residue and the material of generation rather than its whole---the idiom that can be repeated, the rite that can be performed again. \emph{Across the three registers}: the forms are not symbolic only but imaginary, symbolic, and real at once, none dispensable, per the third revision; the culture lives as image and identification, as language and rule, and as the non-symbolizable relational core. \emph{At once generated by the relation and, returning upon it, organizes and sustains it}: the reciprocal turn, in which the exteriorized culture comes back to form the very subjects and shape the very drive that generated it, so that culture is not a product the relation has but a process by which the relation continues to be what it is.

\subsection*{What the definition distinguishes and secures}

Two distinctions the definition draws are worth making explicit, because the rest of the paper turns on them. The first distinguishes culture from the aesthetic field treated in the preceding paper of this series. The aesthetic, there, was the field of present perception, the beauty generated in the relational event of a shared seeing; culture, here, is what that field becomes when it is sedimented over time and made transmissible and re-enactable---the aesthetic field's temporal sedimentation and passage into form. Beauty is generated in the moment of shared perception; culture is generated when such moments deposit into idioms and rites that can be carried forward, so that culture stands to the aesthetic field as the sedimented and repeatable stands to the present and singular. This locates the present paper exactly with respect to its predecessor: where that paper asked how a relation generates present beauty, this one asks how a relation generates a durable world.

The second distinction concerns scale, and it must be stated with the caution the paper has reserved for it. The definition is general; it is meant to hold of culture as such, at any scale. But the paper claims only a \emph{weak analogy}, not a structural identity, between the micro-culture of the intimate dyad and the macro-culture of a people or a symbolic order. The mechanism of emergence---relational subjects exteriorizing interwoven drive and shared history across the registers, under material constraint---may be stated generally, and the general statement illuminates both scales. But the specific features that make the intimate dyad the paper's subject, above all the vacancy of the big Other and the two-person structure that the second part will develop, do not scale up: a people has a full big Other in its traditions and institutions, where a couple has none, and the analogy between the scales must therefore be drawn as an illumination and not pressed into an identity. The definition holds generally; the phenomenology that follows is of the special case, and the special case is special in ways the general definition does not capture. It is to that special case---the emergence of culture in the intimate dyad, in its irreducible particularity---that the paper now turns, and it is there, and not in the general definition just stated, that the paper means to make its distinctive contribution.
